\section{Related work}
\paragraph{HTSR and layer-wise signals.} Martin and Mahoney~\cite{martin2021htsr,martin2021predicting} introduced the power-law fit of layer spectra and the WeightWatcher tool; TempBalance~\cite{zhou2023tempbalance} uses $\alpha$ for layer-wise learning rates on CNNs and AlphaPruning~\cite{lu2024alphapruning} for layer-wise sparsity in LLMs. We transfer the layer-selection idea to the choice of \emph{which block to convert to spiking} and to certifying an insert during fine-tuning. We are not aware of a published $\alpha$ profile of CLIP.

\paragraph{Spiking transformers and conversion.} Spikformer~\cite{zhou2023spikformer} and the Spike-driven Transformer family~\cite{yao2023spikedriven} train spiking ViTs from scratch; SpikeCLIP~\cite{li2023spikeclip} trains a fully spiking CLIP. Conversion methods~\cite{rueckauer2017conversion,you2024spikeziptf,jiang2024sta} convert the whole network and attribute the loss to clipping, quantisation and non-uniform activations. None asks which blocks tolerate conversion or whether a weight statistic predicts it; the common heuristic is to keep softmax and LayerNorm dense and convert MLPs, which is the granularity we study.

\paragraph{Adapters on CLIP.} CLIP-Adapter~\cite{gao2021clipadapter} and Tip-Adapter~\cite{zhang2022tipadapter} add residual adapters or caches to frozen CLIP features for few-shot classification. Our E3 is a CLIP-Adapter with dense, spiking and lookup-memory bodies. Spiking parameter-efficient fine-tuning has so far been studied for spiking backbones (point clouds, neural decoders), not as inserts into a frozen dense vision-language model.
